\section*{अध्याय \ref{ch:b2:oceanic-islands} --- महासागरीय द्वीप}

\begin{solution}{exo:b2:oceanic-islands:1}
(क) पूर्व-शील्ड चरण, आज के लोइही की तरह। (ख) \omterm{def:b2:oceanic-islands:stages}{शील्ड चरण}। (ग) उत्तर-शील्ड चरण। (घ) पुनर्जीवित
चरण। (ङ) अब कोई उद्गारी चरण नहीं: ज्वालामुखी अपनी भित्ति के नीचे धँस गया है, जो अब एक \omterm{def:b2:oceanic-islands:reefs}{एटॉल} है।
\end{solution}

\begin{solution}{exo:b2:oceanic-islands:2}
(क) $2500 + 350\sqrt9 = \qty{3550}{m}$। (ख) $350\times(\sqrt{25} - \sqrt9) = 350\times2 =
\qty{700}{m}$। (ग) $350/(2\times5) = \qty{35}{m}$ प्रति मिलियन वर्ष।
\end{solution}

\begin{solution}{exo:b2:oceanic-islands:3}
(क) मुख्य द्वीपों के आसपास $5000/17 = \qty{294}{kyr}$; कटक के अनुदिश $30\,000/70 = \qty{429}{kyr}$।
(ख) कुछ स्खलन छूट जाते हैं: कोई नया ज्वालामुखी किसी पुराने का निशान और मलबा दबा सकता है,
और पुराने निक्षेप अवसाद से ढक जाते हैं। वास्तविक संख्या अधिक और वास्तविक अंतराल
छोटा है।
\end{solution}

\begin{solution}{exo:b2:oceanic-islands:4}
(क) $4158\times0.3048 = \qty{1267}{m}$; $4610\times0.3048 = \qty{1405}{m}$; $150\times0.3048 =
\qty{46}{m}$। (ख) पहली भित्ति \qty{46}{m} से अधिक गहराई पर नहीं उगी और उसका आधार अब
\qty{1267}{m} पर है: ज्वालामुखी कम से कम $1267 - 46 = \qty{1221}{m}$ धँसा है (अधिक, यदि समुद्र-तल
तब ऊँचा था)। (ग) धँसाव अधिक से अधिक लगभग \qty{1267}{m} है, कम से कम \qty{33.9}{Myr} में:
अधिक से अधिक $1267/33.9 = \qty{37}{m}$ प्रति मिलियन वर्ष, यानी \qty{0.037}{mm} प्रति वर्ष।
\end{solution}

\begin{solution}{exo:b2:oceanic-islands:5}
(क) $\dot S = 175/\sqrt t$: \qty{87.5}{m.Myr^{-1}} (\qty{0.0875}{mm.yr^{-1}}) \qty{4}{Myr} पर,
\qty{35}{m.Myr^{-1}} (\qty{0.035}{mm.yr^{-1}}) \qty{25}{Myr} पर, \qty{25}{m.Myr^{-1}}
(\qty{0.025}{mm.yr^{-1}}) \qty{49}{Myr} पर। (ख) $4.2/0.0875 = 48$ गुना। (ग) $2.6/0.0875 = 30$
गुना: भारण चरण किसी द्वीप को युवा प्लेट के शीतलन से कहीं तेज़ी से धँसाता है।
\end{solution}

\begin{solution}{exo:b2:oceanic-islands:6}
(क) $\dot S = (3200/62.8)\,e^{-100/62.8} = 50.96\times0.2035 = \qty{10.37}{m}$ प्रति मिलियन वर्ष,
यानी प्रति हज़ार वर्ष \qty{0.0104}{m}: मापी गई प्रति हज़ार वर्ष \qty{0.01}{m}। (ख) $e^{-t/62.8}
= 10/50.96$, $t = 62.8\ln5.096 = \qty{102.3}{Myr}$। (ग) मूल नियम: $2500 + 350\times5 =
\qty{4250}{m}$; घातांकी नियम: $6400 - 3200\times0.6716 = \qty{4251}{m}$; उनमें
\qty{1}{m} का अंतर है।
\end{solution}

\begin{solution}{exo:b2:oceanic-islands:7}
(क) $\qty{8.6}{cm.yr^{-1}} = \qty{86}{km.Myr^{-1}}$; $86\times33 = \qty{2838}{km}$। (ख) हाँ:
\qty{27.7}{Myr} उन \qty{33}{Myr} से कम है जिनके बाद भित्तियाँ डूबती हैं, अतः उसकी भित्ति अब भी
जीवित होनी चाहिए। (ग) $55\times1.609 = \qty{88.5}{km}$, यानी $88.5/86 = \qty{1.0}{Myr}$ अधिक: लगभग
\qty{28.7}{Myr}। यह \qty{33}{Myr} से कुछ कम है, पर \qty{33}{Myr} स्वयं पूरी शृंखला के लिए एक
अनुमानित आँकड़ा है, और वृद्धि हर \omterm{def:b2:oceanic-islands:reefs}{एटॉल} पर जल के तापमान पर भी निर्भर करती है:
क्योर, अंतिम जीवित \omterm{def:b2:oceanic-islands:reefs}{एटॉल}, सचमुच \omterm{def:b2:oceanic-islands:darwin-point}{डार्विन बिंदु} के निकट है।
\end{solution}

\begin{solution}{exo:b2:oceanic-islands:8}
(क) बक्सा: $30\times17\times2 = \qty{1020}{km^3}$; पिरामिड: एक तिहाई, \qty{340}{km^3}। (ख) अकेला एक खंड
किलाउआ के पूरे ज्वालामुखी का दसवाँ (बक्सा) या तीसवाँ (पिरामिड) भाग है। (ग) $23\,000\times
0.1 = \qty{2300}{km^3}$, किलाउआ का एक चौथाई; एक ऐसा अनुमान जो मानी गई मोटाई के साथ
ही टिकता या गिरता है।
\end{solution}

\begin{solution}{exo:b2:oceanic-islands:9}
(क) $1000/2.6 = \qty{385}{kyr}$। (ख) हाँ: प्रति हज़ार वर्ष \qty{2.6}{m}
\qty{2.6}{mm.yr^{-1}} है, \qty{4.2}{mm.yr^{-1}} से कम, यानी \qty{1.6}{mm.yr^{-1}} का अंतर।
(ग) समुद्र-तल का सापेक्षिक उठाव $12 + 2.6 = \qty{14.6}{mm.yr^{-1}}$; भित्ति $14.6 - 4.2 =
\qty{10.4}{mm.yr^{-1}}$ से $16.5 - 8.2 = \qty{8.3}{kyr}$ तक पीछे छूटती है: उसका शीर्ष $10.4\times8.3 =
\qty{86}{m}$ पर समाप्त होता है। (घ) $(12 - 4.2)\times8.3 = \qty{65}{m}$। दोनों स्थितियों में वह सबसे गहरी
वृद्धि-सीमा से ऊपर रहती है पर कहीं धुँधले जल में: वह फिर साथ पकड़ पाती है या नहीं, यह इस पर निर्भर है कि गहराई के साथ
उसकी वृद्धि कैसे धीमी होती है।
\end{solution}

\begin{solution}{exo:b2:oceanic-islands:10}
(क) $S(t_1) = C(\sqrt{t_1 + \Delta} - \sqrt{t_1})$ का अवकलज $\frac C2(1/\sqrt{t_1 +
\Delta} - 1/\sqrt{t_1}) < 0$ है: वह $t_1$ के साथ घटता है। (ख) $350\sqrt{20} = \qty{1565}{m}$; $350\times
(\sqrt{70} - \sqrt{50}) = 350\times1.2955 = \qty{453}{m}$। (ग) नहीं: $\sqrt t$ $0$ पर परिमित है, अतः
किसी भी अंतराल में धँसाव परिमित है, (ख) में \qty{1565}{m}। केवल क्षण $t = 0$ पर दर
अनंत है, जहाँ नियम का उद्देश्य अब तल का विस्तार से वर्णन करना नहीं है।
\end{solution}

\begin{solution}{exo:b2:oceanic-islands:11}
(क) $30\times0.1 = \qty{3}{m}$। (ख) लगभग \qty{120}{m}, यदि उच्च-स्तर पर उसका शीर्ष समुद्र-तल पर था:
\omterm{def:b2:oceanic-islands:reefs}{एटॉल} एक चूना-पत्थर का द्वीप बन जाता है, वर्षा प्रवालों के ऐरागोनाइट को घोलती है और
चट्टान गुफाओं के साथ कैल्साइट में पुनःक्रिस्टलित होती है। (ग) $120/3 = 40$ चक्र, \qty{4}{Myr}। (घ) निक्षालित
मंडल अनावरण के यही समय हैं, जो समुद्र लौटने पर फिर उगी भित्ति के नीचे बंद हो गए।
कई चक्रों में समुद्र मिलते-जुलते उच्च-स्तरों पर लौटता है, अतः $\Delta\eta$ छोटा है, भित्ति का शीर्ष
हर उच्च-स्तर पर फिर समुद्र-तल तक बढ़ता है, और टोपी की मोटाई धँसाव से तय होती है, जैसा
प्रस्ताव कहता है।
\end{solution}

\begin{solution}{exo:b2:oceanic-islands:12}
(क) $3200\times(e^{-50/62.8} - e^{-120/62.8}) = 3200\times(0.4510 - 0.1480) = \qty{970}{m}$।
(ख) $150 + 970 = \qty{1120}{m}$। (ग) $700\times1.8288 = \qty{1280}{m}$, अनुमान से \qty{160}{m} अधिक
गहरा। (घ) हो सकता है \omterm{def:b2:oceanic-islands:seamount}{गयो} पहले, किसी युवा प्लेट पर, डूबा हो और अधिक धँसा हो; किसी पड़ोसी
ज्वालामुखी ने उसे भारित किया हो; उसका शीर्ष लहरों ने और नीचे तक काटा हो, या वह अवसाद से ढका हो;
उसके नीचे की प्लेट किसी भिन्न आयु की हो; और डूबने के समय का समुद्र-तल
आज का नहीं था।
\end{solution}

\begin{solution}{pb:b2:oceanic-islands:1}
\textbf{1.} $2500 + 350\sqrt{10} = 2500 + 350\times3.1623 = \qty{3607}{m}$।
\textbf{2.} $3607 + 1000 = \qty{4607}{m}$।
\textbf{3.} $V = \pi r^2h/3 = \pi\times50^2\times4.607/3 = \qty{12061}{km^3}$।
\textbf{4.} $0.85\times12\,061 = \qty{10252}{km^3}$, \qty{1}{Myr} में: प्रति वर्ष \qty{0.0103}{km^3},
आज के किलाउआ से दस गुना कम।
\textbf{5.} $95\,600/12\,061 = 7.9$: लगभग आठ।
\textbf{6.} थोलिआइटी लावा, \omterm{def:b2:oceanic-islands:stages}{शील्ड चरण} से; उत्तर-शील्ड चरण के क्षारीय लावा उसे
ढकते हैं।
\textbf{7.} $\qty{1}{km}/\qty{1}{Myr} = \qty{1}{m}$ प्रति हज़ार वर्ष, हवाई द्वीप के पास की दर से 2.6 गुना
धीमा।
\textbf{8.} हाँ: $4.2 - 1.0 = \qty{3.2}{mm.yr^{-1}}$ अतिरिक्त।
\textbf{9.} $12 + 1 - 4.2 = \qty{8.8}{mm}$ प्रति वर्ष।
\textbf{10.} $150/8.8 = \qty{17}{kyr}$। पार्श्व पर एक डूबी भित्ति-वेदिका, उस प्रकार की सीढ़ी
जो हवाई द्वीप के पास मिलती है।
\textbf{11.} $350\times(\sqrt{22} - \sqrt{10}) = 350\times(4.6904 - 3.1623) = \qty{535}{m}$।
\textbf{12.} $1000 - 535 = \qty{465}{m}$, $465/12 = \qty{38.8}{m}$ प्रति मिलियन वर्ष की दर से।
\textbf{13.} आरंभ में $350/(2\sqrt{10}) = \qty{55.3}{m.Myr^{-1}}$; अंत में $350/(2\sqrt{22}) =
\qty{37.3}{m.Myr^{-1}}$।
\textbf{14.} $2500 + 350\times4.6904 = \qty{4142}{m}$।
\textbf{15.} बड़े हवाई द्वीप लगभग चार गुना तक ऊँचे उठते हैं, अतः समुद्र-तल से ऊपर हटाने के लिए कहीं अधिक
चट्टान है; उनका \omterm{def:b2:oceanic-islands:stages}{शील्ड चरण} भी अधिक समय तक चला।
\textbf{16.} लहरें, जो तट को पीछे काटती हैं और एक तरंग-घर्षित चबूतरा समतल करती हैं; भूस्खलन और \omterm{def:b2:oceanic-islands:collapse}{पार्श्व-ढहन}।
\textbf{17.} एक \omterm{def:b2:oceanic-islands:reefs}{अवरोधक भित्ति}, भित्ति और ज्वालामुखीय तट के अंतिम अंश के बीच एक लैगून के साथ।
\textbf{18.} $350\times(\sqrt{70} - \sqrt{22}) = 350\times(8.3666 - 4.6904) = \qty{1287}{m}$।
\textbf{19.} $1287/48 = \qty{26.8}{m}$ प्रति मिलियन वर्ष, यानी
\qty{0.027}{mm} प्रति वर्ष: भित्ति जितनी बढ़ सकती थी उससे $4.2/0.0268 =
157$ गुना कम।
\textbf{20.} \qty{20}{m} अधिक मोटी, पैरी द्वीप के नीचे बेसाल्ट के ऊपर के चूना-पत्थर से (\qty{1267}{m}),
और \qty{118}{m} पतली, एलुगेलाब के नीचे के चूना-पत्थर से (\qty{1405}{m}): वही परिमाण।
\textbf{21.} $2500 + 350\times8.3666 = \qty{5428}{m}$।
\textbf{22.} $3200\times(e^{-70/62.8} - e^{-120/62.8}) = 3200\times(0.3280 - 0.1480) =
\qty{576}{m}$।
\textbf{23.} $150 + 576 = \qty{726}{m}$।
\textbf{24.} पेट्रेल \omterm{def:b2:oceanic-islands:reefs}{एटॉल} की भित्ति-टोपी \textbf{\qty{1287}{m}} मोटी है, यानी लगभग \qty{1.3}{km}
भित्ति-चूना-पत्थर, जैसा एनीवेटक के नीचे।
\end{solution}
